\section{Introduction}
\label{sec:intro}

\IEEEPARstart{I}{nterpretability} research increasingly communicates through quantitative artifacts. Attribution methods assign relevance to input features~\cite{simonyan2014saliency,sundararajan2017ig,lundberg2017shap}, probes measure what can be decoded from hidden states~\cite{belinkov2022probing}, activation patching measures what changes when an internal value is replaced~\cite{meng2022rome,zhang2024patching}, and circuit and dictionary methods propose components and features as the units of a mechanism~\cite{wang2023ioi,conmy2023acdc,huben2024sae}. Mature libraries make these artifacts cheap to produce~\cite{kokhlikyan2020captum,nanda2022transformerlens,fiottokaufman2025nnsight,wu2024pyvene,bloom2024saelens}.

An artifact, however, is not the claim that is usually attached to it. Consider a BERT-base sentiment classifier and the held-out sentence ``i am sorry that i was unable to get the full brunt of the comedy .'' IG~\cite{sundararajan2017ig} ranks the word \textit{unable} and the separator token \texttt{[SEP]} highest among its 19 positions. That is a statement about one method, one baseline and one target; it does not establish that the two tokens are necessary for the prediction, which is a question about intervention~\cite{jacovi2020faithfulness,joshi2026causality}. When we intervene, the answer depends on choices. Replacing both tokens with the \texttt{[MASK]} embedding lowers the class margin from 3.97 by 2.66, more than half the clean margin and more than the drop caused by at least 95\% of 50 matched random two-token sets; the predicted class does not change. Eight of the 17 declared alternative configurations contradict this result, and if the claim concerns the two highest-ranked \emph{content} tokens the selection becomes $\{$\textit{unable}, \textit{get}$\}$ and the margin drop is negative (Fig.~\ref{fig:example}).

Conclusions of this kind move when the link between measurement and conclusion changes. Activation-patching results depend on the corruption and the metric~\cite{zhang2024patching}; removal-based faithfulness tests depend on how far the replacement is from the data distribution~\cite{hase2021ood,hooker2019roar}; single-input mediation scores have high variance~\cite{meloux2026variance}; subspace interventions can produce illusory localisation~\cite{makelov2024illusion}; probe accuracy does not show that a model uses what the probe finds~\cite{belinkov2022probing,elazar2021amnesic,hewitt2019control}; and agreement with human rationales is not faithfulness~\cite{jacovi2020faithfulness,deyoung2020eraser}. Recent position papers argue that interpretability claims should be stated with their estimand, intervention class and identification assumptions~\cite{joshi2026causality,lin2026identification}, and call for systems that make the arguments behind claims auditable~\cite{lan2026auditable}. What is missing in practice is an executable link: a representation in which the claim, its units, its replacement, its controls and the alternatives that were tried are explicit enough for a program, not only a reviewer, to check the inference.

We present BeyondNN, an open-source evidence and audit layer for PyTorch~\cite{paszke2019pytorch,roufas2026beyondnn}. BeyondNN is not an attribution, probing or patching method. It records each measurement as typed, provenance-bound evidence, represents claims as declared objects with a relation, a subject, a scope and an eligibility, and re-derives each claim's standing from the recorded evidence under a plan that assigns every configuration a role (\textsc{primary}, \textsc{alternative} or \textsc{stress\_test}). Only \textsc{primary} configurations decide a standing; the others produce typed findings. The audit never runs the model and returns no global score.

The contributions of this paper are:
\begin{enumerate}
\item A formal evidence-audit semantics for interpretability claims about neural networks: typed evidence, provenance, scoped claims with explicit eligibility, a test outcome rule with an attainability condition, configuration roles, admissibility and a precedence rule for standings, with six properties that hold by construction (Section~\ref{sec:framework}).
\item An implementation for PyTorch, released as BeyondNN~0.1.0 (Section~\ref{sec:system}).
\item An evaluation on trained MLP, CNN, BERT-tiny and BERT-base models and on compiled and SIIT-trained transformers with reference mechanisms (Sections~\ref{sec:design}--\ref{sec:results}), including: a fixed-selection decomposition of configuration dependence; a comparison with conventional workflows on identical measurements under injected faults; validation against reference labels that use no intervention on the audited network; replication over independently trained checkpoints; exact re-execution; and the computational cost of the evidence layer.
\end{enumerate}

Throughout, we evaluate claim protocols, not attribution methods. IG is the attribution method in these studies; nothing here shows that IG or attribution in general is unreliable. A \textsc{supported} standing means that the evidence passed the declared \textsc{primary} rule for the stated claim, scope, samples and checkpoint; it is not a statement of logical necessity or of a model's mechanism.
